%! Author = lili
%! Date = 4/22/26

\newglossaryentry{frage}{
    name={test},
    description={todo},
    type=frage
}
\newglossaryentry{wasbedingtexpressionsniveaus}{
    name={Was bedingt unterschiedliche Expressions-Niveaus},
    description={todo},
    type=frage
}


\newglossaryentry{wievieleaminoinproteinen}{
    name={Wieviele verschiedene Aminosäuren
    gibt es in Proteinen},
    description={todo},
    type=frage
}

\newglossaryentry{wiefindensichhomologechromosomen}{
    name={Wie finden sich die
    homologen Chromosomen},
    description={todo bw3 F53},
    type=frage
}

\newglossaryentry{warummeioseiinichtmitose}{
    name={Warum ist die Meiose II nicht identisch zur Mitose},
    description={todo},
    type=frage
}

\newglossaryentry{ausnahmeMitose}{
    name={Bei welchen zellen findet keine Verteilung des (replizierten) Chromosomensatzes auf die Tochterzellen statt},
    description={Bei Gameten},
    type=frage
}

\newglossaryentry{nMensch}{
    name={Was ist der Wert n (Chromosomensatz) beim Menschen},
    description={n=23},
    type=frage
}

\newglossaryentry{chromosomenzahlMensch}{
    name={Was ist die Chromosomenzahl des Menschen},
    description={46},
    type=frage
}

\newglossaryentry{washältzsm}{
    name={Was hält Schwesterchromatiden zusammen},
    description={Cohesin-Komplexe},
    type=frage
}

\newglossaryentry{wiechrommetaphase}{
    name={Wie sind die Chromosomen in der Metaphase angeordnet},
    description={todo},
    type=frage
}

\newglossaryentry{wiechromg2}{
    name={Wie sind die Chromosomen in der G2 Phase angeordnet},
    description={todo},
    type=frage
}

\newglossaryentry{wiechromg1}{
    name={Wie sind die Chromosomen in der G1 Phase angeordnet},
    description={todo},
    type=frage
}

\newglossaryentry{wasinmitose}{
    name={Was passiert in der Mitose},
    description={Sie Chromosomen ordnen sich in der Metaphasenplatte, Schwesterchromatiden trennen sich und es entstehen zwei Tochterzellen mit je 2n. Anders gesagt: Verteilung des (replizierten) Chromosomensatzes bei der
    Vermehrung von (somatischen) Zellen auf die Tochterzellen},
    type=frage
}




\newglossaryentry{wannchromo}{
    name={Wann sind Chromosomen zu sehen},
    description={in der Mitose},
    type=frage
}

\newglossaryentry{wenigac}{
    name={Was bedeutet es, wenn Nukleosomen wenig Ac haben},
    description={todo},
    type=frage
}

\newglossaryentry{dnatrennung}{
    name={Wie trennt man doppelsträngigge DNA auf},
    description={todo},
    type=frage
}


\newglossaryentry{avery}{
    name={Wie hat Avery herausgefunden, dass die Nukleinsäure der Träger der Erbinformationen ist},
    description={todo},
    type=frage
}

\newglossaryentry{menschchromosomen}{
    name={Wie viele Chromosomen welcher Art hat ein Mensch},
    description={46 Chromosomen: 2 x 22 Autosomen, 2 Geschlechtschromosomen},
    type=frage
}